\section{Homework Week 9}

\begin{exercise}
    Let $\varphi$ be a multiplicative valuation of a field $K$. Then the following hold:
    \begin{enumerate}
        \item The sequence $(a_n)$ is a Cauchy sequence if and only if
        \[
            \lim_{n\to\infty}\varphi(a_{n+1}-a_n)=0.
        \]
    
        \item If
        \[
            \lim_{n\to\infty}a_n=a\neq 0,
        \]
        then
        \[
            \varphi(a_n)=\varphi(a)
        \]
        for all sufficiently large $n$.
    \end{enumerate}
    \end{exercise}

    \begin{proof}
        We use the metric on $K$ induced by $\varphi$, namely
        \[
            d(x,y)=\varphi(x-y).
        \]
        
        \begin{enumerate}
            \item Suppose first that $(a_n)$ is a Cauchy sequence. Then, for every $\varepsilon>0$,
            there exists $N$ such that for all $m,n\geq N$,
            \[
                \varphi(a_m-a_n)<\varepsilon.
            \]
            Taking $m=n+1$, we get
            \[
                \varphi(a_{n+1}-a_n)<\varepsilon
            \]
            for all $n\geq N$. Hence
            \[
                \lim_{n\to\infty}\varphi(a_{n+1}-a_n)=0.
            \]
        
            Conversely, suppose that
            \[
                \lim_{n\to\infty}\varphi(a_{n+1}-a_n)=0.
            \]
            We show that $(a_n)$ is Cauchy. Let $\varepsilon>0$ be given. Then there exists
            $N$ such that for all $k\geq N$,
            \[
                \varphi(a_{k+1}-a_k)<\varepsilon.
            \]
            If $m>n\geq N$, then
            \[
                a_m-a_n
                =
                (a_m-a_{m-1})+(a_{m-1}-a_{m-2})+\cdots+(a_{n+1}-a_n).
            \]
            Since $\varphi$ is non-archimedean, we have
            \[
                \varphi(a_m-a_n)
                \leq
                \max_{n\leq k\leq m-1}\varphi(a_{k+1}-a_k)
                <
                \varepsilon.
            \]
            Therefore $(a_n)$ is a Cauchy sequence.
        
            \item Suppose that
            \[
                \lim_{n\to\infty}a_n=a\neq 0.
            \]
            Then
            \[
                \lim_{n\to\infty}\varphi(a_n-a)=0.
            \]
            Since $a\neq 0$, we have $\varphi(a)>0$. Hence there exists $N$ such that for all
            $n\geq N$,
            \[
                \varphi(a_n-a)<\varphi(a).
            \]
            For such $n$, we write
            \[
                a_n=a+(a_n-a).
            \]
            By the strong triangle inequality,
            \[
                \varphi(a_n)
                \leq
                \max\{\varphi(a),\varphi(a_n-a)\}
                =
                \varphi(a).
            \]
            On the other hand,
            \[
                a=a_n-(a_n-a),
            \]
            so
            \[
                \varphi(a)
                \leq
                \max\{\varphi(a_n),\varphi(a_n-a)\}.
            \]
            Since $\varphi(a_n-a)<\varphi(a)$, the maximum on the right-hand side must be
            $\varphi(a_n)$. Thus
            \[
                \varphi(a)\leq \varphi(a_n).
            \]
            Combining the two inequalities gives
            \[
                \varphi(a_n)=\varphi(a)
            \]
            for all $n\geq N$.
        \end{enumerate}
        \end{proof}
    
    \begin{exercise}
    Let $\varphi$ be a multiplicative valuation of a field $K$. Let $V$ be a finite dimensional
    $K$-space with basis $u_1,\dots,u_n$. For
    \[
        v=a_1u_1+\cdots+a_nu_n\in V,
    \]
    define $\|v\|_0$ by
    \[
        \|v\|_0=\max\{\varphi(a_i)\mid 1\leq i\leq n\}.
    \]
    Then this is a norm on $V$. Moreover, the following hold concerning the metric induced by
    this norm:
    \begin{enumerate}
        \item Let
        \[
            v_r=a_{r1}u_1+\cdots+a_{rn}u_n\in V
        \]
        be given for $r=1,2,\dots$. Then the sequence $(v_r)$ is a Cauchy sequence if and only if
        the sequence $(a_{ri})_r$ in $K$ is a Cauchy sequence for every $i$.
    
        Also,
        \[
            \lim_{r\to\infty}v_r=\sum_{i=1}^n a_i u_i
        \]
        if and only if
        \[
            \lim_{r\to\infty}a_{ri}=a_i
        \]
        for every $i$.
    
        \item If $K$ is complete, then so is $V$.
    \end{enumerate}
    \end{exercise}

    \begin{proof}
        First we prove that $\|\cdot\|_0$ is a norm on $V$.
        
        Let
        \[
            v=a_1u_1+\cdots+a_nu_n.
        \]
        Then
        \[
            \|v\|_0=\max\{\varphi(a_i)\mid 1\leq i\leq n\}.
        \]
        Since $\varphi(a_i)\geq 0$ for every $i$, we have
        \[
            \|v\|_0\geq 0.
        \]
        Moreover,
        \[
            \|v\|_0=0
        \]
        if and only if
        \[
            \varphi(a_i)=0
        \]
        for every $i$, which is equivalent to
        \[
            a_i=0
        \]
        for every $i$. Since $u_1,\dots,u_n$ is a basis, this is equivalent to
        \[
            v=0.
        \]
        
        Now let $\lambda\in K$. Then
        \[
            \lambda v
            =
            \lambda a_1u_1+\cdots+\lambda a_nu_n,
        \]
        and therefore
        \[
            \|\lambda v\|_0
            =
            \max_{1\leq i\leq n}\varphi(\lambda a_i)
            =
            \max_{1\leq i\leq n}\varphi(\lambda)\varphi(a_i).
        \]
        Hence
        \[
            \|\lambda v\|_0
            =
            \varphi(\lambda)\max_{1\leq i\leq n}\varphi(a_i)
            =
            \varphi(\lambda)\|v\|_0.
        \]
        
        Finally, let
        \[
            v=a_1u_1+\cdots+a_nu_n,
            \qquad
            w=b_1u_1+\cdots+b_nu_n.
        \]
        Then
        \[
            v+w=(a_1+b_1)u_1+\cdots+(a_n+b_n)u_n.
        \]
        Thus
        \[
            \|v+w\|_0
            =
            \max_{1\leq i\leq n}\varphi(a_i+b_i).
        \]
        Since $\varphi$ is non-archimedean,
        \[
            \varphi(a_i+b_i)\leq \max\{\varphi(a_i),\varphi(b_i)\}
        \]
        for every $i$. Hence
        \[
            \|v+w\|_0
            \leq
            \max_{1\leq i\leq n}\max\{\varphi(a_i),\varphi(b_i)\}.
        \]
        Therefore
        \[
            \|v+w\|_0
            \leq
            \max\left\{
                \max_{1\leq i\leq n}\varphi(a_i),
                \max_{1\leq i\leq n}\varphi(b_i)
            \right\}
            =
            \max\{\|v\|_0,\|w\|_0\}.
        \]
        So $\|\cdot\|_0$ is a non-archimedean norm on $V$.
        
        \begin{enumerate}
            \item Let
            \[
                v_r=a_{r1}u_1+\cdots+a_{rn}u_n.
            \]
            For $r,s\geq 1$, we have
            \[
                v_r-v_s
                =
                (a_{r1}-a_{s1})u_1+\cdots+(a_{rn}-a_{sn})u_n.
            \]
            Hence
            \[
                \|v_r-v_s\|_0
                =
                \max_{1\leq i\leq n}\varphi(a_{ri}-a_{si}).
            \]
        
            Suppose first that $(v_r)$ is Cauchy. Let $i$ be fixed. Since
            \[
                \varphi(a_{ri}-a_{si})
                \leq
                \max_{1\leq j\leq n}\varphi(a_{rj}-a_{sj})
                =
                \|v_r-v_s\|_0,
            \]
            it follows that $(a_{ri})_r$ is a Cauchy sequence in $K$.
        
            Conversely, suppose that for every $i$, the sequence $(a_{ri})_r$ is Cauchy in
            $K$. Let $\varepsilon>0$ be given. For each $i$, there exists $N_i$ such that for
            all $r,s\geq N_i$,
            \[
                \varphi(a_{ri}-a_{si})<\varepsilon.
            \]
            Let
            \[
                N=\max\{N_1,\dots,N_n\}.
            \]
            Then for all $r,s\geq N$ and all $i$, we have
            \[
                \varphi(a_{ri}-a_{si})<\varepsilon.
            \]
            Therefore
            \[
                \|v_r-v_s\|_0
                =
                \max_{1\leq i\leq n}\varphi(a_{ri}-a_{si})
                <
                \varepsilon.
            \]
            Hence $(v_r)$ is Cauchy in $V$.
        
            Next we prove the assertion about convergence. Let
            \[
                v=\sum_{i=1}^n a_i u_i.
            \]
            Then
            \[
                v_r-v
                =
                (a_{r1}-a_1)u_1+\cdots+(a_{rn}-a_n)u_n.
            \]
            Hence
            \[
                \|v_r-v\|_0
                =
                \max_{1\leq i\leq n}\varphi(a_{ri}-a_i).
            \]
            Therefore
            \[
                \lim_{r\to\infty}v_r=v
            \]
            if and only if
            \[
                \lim_{r\to\infty}\max_{1\leq i\leq n}\varphi(a_{ri}-a_i)=0.
            \]
            Since there are only finitely many indices $i$, this is equivalent to
            \[
                \lim_{r\to\infty}\varphi(a_{ri}-a_i)=0
            \]
            for every $i$, which is precisely
            \[
                \lim_{r\to\infty}a_{ri}=a_i
            \]
            for every $i$.
        
            \item Suppose that $K$ is complete. We show that $V$ is complete.
        
            Let $(v_r)$ be a Cauchy sequence in $V$, where
            \[
                v_r=a_{r1}u_1+\cdots+a_{rn}u_n.
            \]
            By part $(1)$, for each $i$, the sequence $(a_{ri})_r$ is Cauchy in $K$.
            Since $K$ is complete, there exists $a_i\in K$ such that
            \[
                \lim_{r\to\infty}a_{ri}=a_i.
            \]
            Define
            \[
                v=\sum_{i=1}^n a_i u_i\in V.
            \]
            Again by part $(1)$, we get
            \[
                \lim_{r\to\infty}v_r=v.
            \]
            Thus every Cauchy sequence in $V$ converges in $V$. Hence $V$ is complete.
        \end{enumerate}
        \end{proof}